\documentclass[10pt,a4paper]{amsart}
\usepackage[margin=19mm]{geometry}
\usepackage{amsmath,amssymb,mathtools}
\usepackage[scr=rsfs,cal=euler]{mathalfa}
\usepackage{xfrac}
\usepackage[unicode,hidelinks,bookmarks=false]{hyperref}
\newcommand{\ie}{i.e.\ }
\newcommand{\cf}{cf.\ }
\newcommand{\resp}{respectively\ }
\newcommand{\Subsection}{\subsection}
\providecommand{\nobreakdash}{\nobreak\hbox{-}}
\setlength{\emergencystretch}{2em}
\multlinegap0pt
\usepackage{amssymb}
\usepackage{bm}
\usepackage{enumerate}
\newtheorem{theorem}{Theorem}
\title{Weak-Strong Uniqueness for a Rigid Body Immersed in an Inviscid Compressible Fluid}
\author{Qianfeng Li, Emil Wiedemann}
\date{Source: arXiv:2603.03151v1 (CC BY 4.0)}
\begin{document}
\maketitle

\noindent\emph{Source and licence.} Weak-Strong Uniqueness for a Rigid Body Immersed in an Inviscid Compressible Fluid. Version arXiv:2603.03151v1 (CC BY 4.0), distributed by arXiv under the Creative Commons Attribution 4.0 International licence, http://creativecommons.org/licenses/by/4.0/, as recorded in the arXiv licence metadata for this exact version and shown on the abstract page https://arxiv.org/abs/2603.03151v1. Full licence notice: \texttt{licenses/arxiv-2603.03151-CC-BY-4.0.txt}.\par\smallskip

\noindent\textit{Editorial scope.} Counted unit: the article's theorem thm:ApendixA (source0.tex lines 1834--1854). The article's complete original proof of it is delivered with it (source0.tex lines 1856--1964, one \texttt{proof} environment of 109 lines). No mathematics is written, repaired or re-derived: the statement and the proof are the article's own lines, in the article's own notation, and no retained line is substituted or re-flowed.  No further source-local context is needed: every cross-reference of every retained line resolves inside the two delivered spans.  Nothing was omitted from any retained span, so the package carries no omission record.  No retained line contains a citation, so the wrapper carries no bibliography and no bibliography entry.  No retained line contains a figure environment, an image-inclusion command or drawing code, so no image is delivered and the package declares no asset dependency.  Editorial formatting only, in addition: an AMS \texttt{amsart} preamble that restates the article's own declarations and macros used by the retained lines, namely the environments \texttt{theorem} on separate counters, together with the standard packages those lines need.  The retained environments print in this document as Theorem 1 in order of appearance, whereas the article prints the same items with the numbering of its own full document.  The complete native source of the article as distributed in the arXiv e-print for this exact version is bundled unchanged as \texttt{sources/195657/source0.tex}.
\begin{theorem}\label{thm:ApendixA}[From Newton--Euler to the distributional momentum equation with a surface delta]

In the inviscid compressible fluid-body system, suppose the Newton--Euler balances hold for all $t>0$:
\begin{align}
M\,{V}'(t) &= \int_{\partial\mathcal{B}_t}  p_Fn \, \mathrm{d}S, \label{NE-lin}\\
\frac{d}{dt}\big(\mathcal J(t)\,\omega(t)\big)
&= \int_{\partial\mathcal{B}_t} r \times p_Fn \, \mathrm{d}S, \label{NE-ang}
\end{align}
Here and in the sequel $r=x-X(t)$ and $X(t)$ is the mass center. Then the following \emph{distributional} momentum equation for the rigid phase holds on $\mathbb R^3\times(0,T)$:
\begin{equation}
\partial_t\big(\chi_{\mathcal{B}_t}\rho_B\, v\big)+\nabla\!\cdot\!\big(\chi_{\mathcal{B}_t}\rho_B\, v\otimes v\big)
= \big( p_Fn\big)\,\delta_{\partial\mathcal{B}_t}
\label{PDE-delta}
\end{equation}
for test functions in $C_c^\infty((0,T);\mathcal{R}(\Omega))$ where  
\begin{equation}
    \mathcal{R}(\Omega):=\{\xi:\Omega\to\mathbb{R}^3 \big| \xi=\alpha+\eta\times(\cdot-X), \alpha,\eta,X\in \mathbb{R}^3 \},
\end{equation}
$\chi_{\mathcal{B}_t}$ is the indicator of $\mathcal B_t$, and $\delta_{\partial\mathcal{B}_t}$ is the surface delta distribution
characterized by $\int_{\mathbb R^3}\phi\,\delta_{\partial\mathcal{B}_t}\,\mathrm{d}x=\int_{\partial\mathcal{B}_t}\phi\,\mathrm{d}S$.
\end{theorem}

\begin{proof}
We show that \eqref{PDE-delta} holds in the sense of distributions by verifying its space-time weak
form. Let $w\in C_c^\infty([0,T];\mathcal{R}(\Omega))$ be an arbitrary test function with
$w(\cdot,T)=0$. We need to prove
\begin{equation}
-\!\!\int_0^T\!\!\int_{\mathbb R^3} \chi_{\mathcal{B}_t}\rho_B\,v\cdot\partial_tw\,\mathrm{d}x\,\mathrm{d}t
-\!\!\int_0^T\!\!\int_{\mathbb R^3} \chi_{\mathcal{B}_t}\rho_B\,v\otimes v:\nabla w\,\mathrm{d}x\,\mathrm{d}t
=\int_0^T\!\!\int_{\partial\mathcal{B}_t} p_Fn\cdot w\,\mathrm{d}S\,\mathrm{d}t.
\label{weak}
\end{equation}
Because $\chi_{\mathcal{B}_t}$ localizes the integrals to $\mathcal B_t$, the left-hand side equals
\begin{equation}
-\!\!\int_0^T\!\!\int_{\mathcal B_t} \rho_B\,v\cdot\partial_t w\,\mathrm{d}x\,\mathrm{d}t
-\!\!\int_0^T\!\!\int_{\mathcal B_t} \rho_B\, v\otimes v:\nabla w\,\mathrm{d}x\,\mathrm{d}t.
\label{LHS}
\end{equation}

\noindent\textbf{Step 1: A kinematic transport identity.}
On the moving domain $\mathcal B_t$ with boundary velocity $v$, Reynolds' transport theorem
and the product rule yield, for every smooth $w$,
\begin{equation}
\frac{\mathrm{d}}{\mathrm{d}t}\int_{\mathcal B_t}\rho_B\,v\cdot w\,\mathrm{d}x
=\int_{\mathcal B_t}\rho_B\, v\cdot\partial_t w\,\mathrm{d}x
+\int_{\mathcal B_t}\rho_B\, v\otimes v:\nabla w\,\mathrm{d}x
+\int_{\mathcal B_t}\rho_B\, a\cdot w\,\mathrm{d}x,
\label{K0}
\end{equation}
where $ a:=\partial_t v+(v\!\cdot\!\nabla) v$ is the Eulerian acceleration.
Identity \eqref{K0} is purely kinematic (it uses mass conservation in the rigid body:
$\partial_t\rho_B+\operatorname{div}(\rho_B v)=0$).

\noindent\textbf{Step 2: Pairing the acceleration with the test function.}
For an arbitrary $w\in C_c^\infty([0,T];\mathcal{R}(\Omega))$, there exists $\alpha(t), \eta(t)$ such that 
\begin{equation}\label{proj}
   w(x,t)=\alpha(t)+\eta(t)\times r,\quad r=x-X(t). 
\end{equation}



Then the following pairing identity holds:
\begin{equation}
\int_{\mathcal B_t}\rho_B\,a\cdot w\,\mathrm{d}x
= M\,{V'}\cdot\alpha + \frac{\mathrm{d}}{\mathrm{d}t}\big(\mathcal J\omega\big)\cdot\eta.
\label{pair}
\end{equation}
Indeed, using the rigid kinematics
$a={V'}+{\omega'}\times r+\omega\times(\omega\times r)$,
the first term integrates to $M{V'}\cdot \alpha$ by \eqref{proj}, and the remaining
terms give, by the scalar triple-product identity and the identity
\begin{equation}
\frac{\mathrm{d}}{\mathrm{d}t}\big(\mathcal J\omega\big)
=\int_{\mathcal B_t}\rho_B\,\big(r\times a\big)\,\mathrm{d}x,
\label{Ldot}
\end{equation}
the second contribution in \eqref{pair}. Equality \eqref{Ldot} is obtained by differentiating the
angular momentum $\int_{\mathcal B_t}\rho_B\,( r\times v)\,\mathrm{d}x=\mathcal J\omega$
with Reynolds' theorem and using $\int_{\mathcal B_t}\rho_B\,(V\times v)\,\mathrm{d}x
= 0$.

\noindent\textbf{Step 3: Insert \eqref{pair} into \eqref{K0} and integrate in time.}
Combining \eqref{K0} and \eqref{pair} and rearranging gives
\begin{equation}
-\!\!\int_{\mathcal B_t}\rho_B\, v\cdot\partial_t w\,\mathrm{d}x
-\!\!\int_{\mathcal B_t}\rho_B\, v\otimes v:\nabla w\,\mathrm{d}x
= M\,{V'}\cdot \alpha + \frac{\mathrm{d}}{\mathrm{d}t}\big(\mathcal J\omega\big)\cdot\eta
-\frac{\mathrm{d}}{\mathrm{d}t}\int_{\mathcal B_t}\rho_B\,v\cdot w\,\mathrm{d}x.
\label{K}
\end{equation}
Integrating \eqref{K} over $(0,T)$ and using $w(\cdot,0)=w(\cdot,T) =0$ yields
%\begin{align}
%&-\!\!\int_0^T\!\!\int_{\mathcal B_t}\rho_B\, v\cdot\partial_t w\,\mathrm{d}x\,\mathrm{d}t
%-\!\!\int_0^T\!\!\int_{\mathcal B_t}\rho_B\, v\otimes v:\nabla w\,\mathrm{d}x\,\mathrm{d}t \nonumber\\
%&\qquad= \int_0^T \Big(M\,{V'}\cdot \alpha
%+\frac{\mathrm{d}}{\mathrm{d}t}\big(\mathcal J\omega\big)\cdot\eta\Big)\,\mathrm{d}t
%+\int_{\mathcal B_0}\rho_B\, v(\cdot,0)\cdot w(\cdot,0)\,\mathrm{d}x.
%\label{pre-weak}
%\end{align}
%The initial term can be absorbed in the usual way by requiring $w(\cdot,0)= 0$ (or by
%recording initial data separately). With $w(\cdot,0)= 0$, we obtain
\begin{equation}
-\!\!\int_0^T\!\!\int_{\mathcal B_t}\rho_B\,v\cdot\partial_tw\,\mathrm{d}x\,\mathrm{d}t
-\!\!\int_0^T\!\!\int_{\mathcal B_t}\rho_B\,v\otimes v:\nabla w\,\mathrm{d}x\,\mathrm{d}t
= \int_0^T \Big(M\,{V'}\cdot \alpha
+\frac{\mathrm{d}}{\mathrm{d}t}\big(\mathcal J\omega\big)\cdot\eta\Big)\,\mathrm{d}t.
\label{weak-left}
\end{equation}

\noindent\textbf{Step 4: Invoke Newton--Euler and identify the interface work.}
By the Newton--Euler relations \eqref{NE-lin}--\eqref{NE-ang},
\[
M\,{V'}\cdot \alpha
+\frac{\mathrm{d}}{\mathrm{d}t}\big(\mathcal J\omega\big)\cdot\eta
=\int_{\partial\mathcal{B}_t}p_Fn\cdot
\big(\alpha+\eta\times r\big)\,\mathrm{d}S=\int_{\partial\mathcal{B}_t} p_Fn\cdot w\,\mathrm{d}S.
\]
Therefore,
\[
\int_0^T \Big(M\,{V'}\cdot \alpha
+\frac{\mathrm{d}}{\mathrm{d}t}\big(\mathcal J\omega\big)\cdot\eta\Big)\,\mathrm{d}t
=\int_0^T\!\!\int_{\partial\mathcal{B}_t} p_Fn\cdot w\,\mathrm{d}S\,\mathrm{d}t.
\]
Substituting this into \eqref{weak-left} gives \eqref{weak}. Thus the distribution
identity \eqref{PDE-delta} follows by the definition of the surface delta:
\[
\int_0^T\!\!\int_{\mathbb R^3}\big(p_Fn\,\delta_{\partial\mathcal{B}_t}\big)\cdot w\,\mathrm{d}x\,\mathrm{d}t
=\int_0^T\!\!\int_{\partial\mathcal{B}_t} p_Fn\cdot w\,\mathrm{d}S\,\mathrm{d}t.
\]
This completes the proof.
\end{proof}

\end{document}
